%% Elements of Nested Fibrational Cosmology
%% LING Branch Book: Linguistics Branch (Prospective)
%% Status: Architecturally constituted; pre-closure — all theorem shells carry open obligations
%%
%% Status legend:
%%   [U]  Unconditional  -- proved from spine imports and prior [U] results
%%   [C]  Conditional    -- depends on explicitly declared hypotheses
%%   [D]  Definition-structural
%%   [R]  Remark only
%%   [O]  Open obligation / named failure frontier
%%   [B]  Bridge (Book VI licensed continuum surrogate)

\documentclass[12pt,oneside]{amsbook}
\usepackage{amsmath,amssymb,amsthm}
\usepackage{mathrsfs}
\usepackage{enumitem}
\usepackage{hyperref}
\usepackage{geometry}
\geometry{margin=1.2in}

\theoremstyle{plain}
\newtheorem{theorem}{Theorem}[section]
\newtheorem{lemma}[theorem]{Lemma}
\newtheorem{proposition}[theorem]{Proposition}
\newtheorem{corollary}[theorem]{Corollary}

\theoremstyle{definition}
\newtheorem{definition}[theorem]{Definition}
\newtheorem{standingrule}[theorem]{Standing Rule}
\newtheorem{openobligation}[theorem]{Open Obligation}

\theoremstyle{remark}
\newtheorem{remark}[theorem]{Remark}
\newtheorem{scholium}[theorem]{Scholium}

\newcommand{\status}[1]{\textup{[#1]}}

%% Notation
%%   R_Ling              Linguistics target regime
%%   O_Ling              Linguistics observable family = (K_Ling, C_Ling, Gamma_Ling)
%%   K_Ling              Contrast kernel (quotient-visible distinguishability data)
%%   C_Ling              Compositional ledger (hierarchical assembly records)
%%   Gamma_Ling          Grammar-descent object (certified rule object)
%%   Delta_Ling          Contrast margin (minimum separation of equivalence classes)
%%   E_Ling              Linguistics endpoint datum
%%   Resp_L(X;P)         Quotient-visible contrast signature under admissible probe P

\begin{document}

\title{LING Branch Book\\[6pt]
       Contrast, Compositionality, and Certified Linguistic Structure}
\author{Elements of Nested Fibrational Cosmology\\
        Prospective Derived Branch}
\date{}
\maketitle
\tableofcontents
\newpage

%% ---------------------------------------------------------------
\section{Purpose and Canonical Seed}\label{sec:ling-purpose}
%% ---------------------------------------------------------------

This branch studies whether linguistic structure --- phonological
contrast, syntactic compositionality, semantic equivalence --- can be
realized as a lawful descendant of the NFC source architecture.  Its
task is not to import the standard language of phonemes, grammars,
lexicons, or meaning-bearing expressions as primitives.  Its task is to
determine when such language is licensed as an effective encoding of
certified discrete observable structure within a declared branch regime.
That posture is required by Book~I's no-smuggling stance, Book~V's
branch-legitimacy law, Book~VI's continuum-interface discipline, and
Book~VII's scoped-certificate law.

The central question is: starting from nothing but finite signal-object
pairs, a finite admissible contrast discipline, and the observational
equivalence imposed by that discipline, can one recover certified
hierarchical compositional structure?  No alphabet is assumed.  No
phoneme inventory is assumed.  No syntactic category system is assumed.
No semantic domain is assumed.  Each such notion, if it appears at all,
must emerge as a quotient-visible, source-descended object.

The linguistics branch is motivated by two structural templates already
present in the canonical corpus.  First, the SPEC branch already treats
probe-response pairs and transition ledgers as the building block of a
certified observable family --- the \emph{contrast-probe template}.
Second, the NS branch already treats a continuation criterion and
obstruction-decay pattern as the backbone of its main frontier theorem
--- a template for \emph{compositional propagation with named failure
frontier}.  A linguistics branch unifies these two tendencies: contrast
structure (analogous to spectral data) plus compositional assembly
(analogous to one-step interface propagation).

The initial endpoint is deliberately modest: certified contrast
structure and one-step compositional assembly on a declared linguistic
target regime.  Full semantic compositionality, infinite generativity
(recursion), and cross-linguistic universals are deferred as named open
obligations until the foundational contrast and compositional layers are
genuinely discharged.

\begin{remark}[\status{R}]\label{rem:ling-strategic-path}
\textbf{Strategic path.}
The shortest canonical route is contrast-founded compositionality:
begin from certified contrast equivalence classes (the discrete
phonological layer), assemble a one-step compositional ledger from
admissible probe episodes, prove a contrast-margin theorem, then prove a
no-hidden-channel theorem, then ask whether the assembled structure
certifiably descends from $\mathbf{U}$ via the Book~V branch-legitimacy
law.  Semantic grounding and recursion-licensing remain downstream open
obligations.
\end{remark}

%% ---------------------------------------------------------------
\section{Imported Spine Structure}\label{sec:ling-imports}
%% ---------------------------------------------------------------

\textbf{Imported spine results.}  All certified results of Books~I--VII
are required for lawful descent, transport accounting, continuum-interface
discipline, and closure governance.  In particular:
\begin{enumerate}[label=\textup{(\arabic*)}]
  \item \emph{Book~I}: observational quotient discipline; no primitive
    symbols, alphabets, syntactic categories, or hidden naming conventions
    may be assumed.  Linguistic content must therefore be
    quotient-visible and admissibly testable, not imported as background
    linguistics.
  \item \emph{Book~II}: stabilized local outcomes, collar-visible
    control, and boundary capacity structure.  The finiteness of the
    observable contrast boundary governs what can be read off from a
    finite signal-object boundary.
  \item \emph{Book~III}: persistence and transport law for certified
    observables; additive bookkeeping.  Compositional transport costs
    must be accountable within this framework.
  \item \emph{Book~IV}: the universal source object $\mathbf{U}$ and
    source-descended comparison architecture.  All certified linguistic
    objects must factor through $\mathbf{U}$.
  \item \emph{Book~V}: lawful branchhood and branch visibility; the
    constitutional law under which a linguistics branch is admitted.
  \item \emph{Book~VI}: continuum language licensed only as faithful
    shorthand for certified discrete structure within a declared interface
    regime; the main gate for acoustic frequency notation, probability
    distributions over utterances, and information-theoretic quantities
    in formal linguistic models.
  \item \emph{Book~VII}: scoped certificates, promotion law, named
    frontier stability, and prohibition of hidden upgrade.  Claims about
    universal grammar, infinite generativity, or cross-linguistic
    universals are subject to this governance without exception.
\end{enumerate}

\textbf{Imported branch precedents.}
\begin{enumerate}[label=\textup{(\arabic*)}]
  \item \emph{SPEC branch}: probe-response objects and transition
    ledgers as building blocks of a certified observable family; the
    \emph{contrast-probe template} for this branch.
  \item \emph{NS branch}: continuation criterion and obstruction-decay
    pattern as the backbone of a main frontier theorem; the
    \emph{compositional propagation with named failure frontier} template.
  \item \emph{SM and matter branches}: category and symmetry content
    remain downstream burdens rather than primitive branch assumptions;
    no linguistic universal is to be treated as background category
    theory.
\end{enumerate}

%% ---------------------------------------------------------------
\section{Branch-Specific Arena}\label{sec:ling-arena}
%% ---------------------------------------------------------------

\subsection{Standing Rules}\label{sec:ling-standing-rules}

\begin{standingrule}[\status{D}]\label{sr:ling-no-smuggling}
No linguistic object --- phoneme, morpheme, word, phrase, sentence,
proposition, or meaning --- may be treated as primitive branch data
merely because it is familiar in external linguistic theory.  Every
branch-visible linguistic object must be sourced through $\mathbf{U}$,
the declared descent machinery, and the linguistics observable family.
This is a branch-local instance of Book~I no-smuggling and Book~V
branch-legitimacy discipline.
\end{standingrule}

\begin{standingrule}[\status{D}]\label{sr:ling-quotient-visible}
A linguistic distinction counts only when it is quotient-visible under
admissible contrast tests or under a lawfully certified contrast
signature descended from such tests.  Book~I's observational quotient
discipline remains controlling.  In particular, a phonological
distinction between two signal-objects counts only when it is
witnessable by some admissible contrast probe.
\end{standingrule}

\begin{standingrule}[\status{D}]\label{sr:ling-no-hidden-symbol}
No hidden alphabet, feature matrix, underlying representation, or
abstract phonological level may be assumed.  Symbol-level structure may
appear only after a branch-local derivation certifies that such
structure faithfully encodes quotient-visible contrast data on the
declared regime.
\end{standingrule}

\begin{standingrule}[\status{D}]\label{sr:ling-no-self-promotion}
A finite contrast certificate may not self-promote to a claim about
infinite generativity, recursion, or universal grammar.  Such promotion
requires an explicit transfer theorem.  This is a direct Book~VII
governance consequence.
\end{standingrule}

\begin{standingrule}[\status{D}]\label{sr:ling-frontier-visibility}
Named frontiers in compositional structure, recursive generativity,
semantic grounding, or cross-linguistic universals remain visible until
discharged by theorem.  Rephrasing does not discharge them.
\end{standingrule}

\begin{standingrule}[\status{D}]\label{sr:ling-no-semantic-primitive}
No semantic domain, truth-value space, or model-theoretic universe may
be treated as a primitive.  Meaning, to the extent that it appears at
all, must emerge as a quotient-visible assignment licensed by certified
contrast structure.
\end{standingrule}

\begin{standingrule}[\status{D}]\label{sr:ling-continuum-gate}
Continuous acoustic quantities, probability distributions over
utterances, and information-theoretic rate quantities may appear only
after a branch-local Book~VI bridge theorem certifies that they
faithfully encode certified discrete contrast structure within a
declared continuum interface regime.
\end{standingrule}

\subsection{Declared Target Regime}\label{sec:ling-target-regime}

\begin{definition}[\status{D}]\label{def:ling-target-regime}
\textup{(Linguistics Target Regime.)}
A \emph{linguistics target regime} is a declared branch regime
$\mathcal{R}_{\mathrm{Ling}}$ consisting of:
\begin{enumerate}[label=\textup{(\arabic*)}]
  \item a certified observable class $\mathcal{C}_{\mathrm{obs}}$ of
    signal-object pairs --- the objects on which contrast tests are
    performed;
  \item a finite admissible contrast discipline $\mathcal{P}$, each
    element of which is a declared admissible process detecting a
    quotient-visible distinction between signal-object pairs;
  \item a declared compositional window $\mathcal{W}_{\mathrm{comp}}$
    specifying the range of assembly operations within which
    compositional claims have force; and
  \item a declared comparison discipline governing when two assembled
    objects are equivalent for the purpose of contrast and
    compositional claims.
\end{enumerate}
The role of $\mathcal{R}_{\mathrm{Ling}}$ is to state explicitly where
linguistic claims are supposed to have force.
\end{definition}

\begin{remark}[\status{R}]\label{rem:ling-signal-object}
A \emph{signal-object} is not yet a phonological string, a word, or a
sentence.  At this stage it is only a certified observable configuration
on which admissible contrast tests can be performed.  The quotient of
the signal-object space under admissible contrast equivalence will, if
the branch succeeds, yield the certified contrast class structure
analogous to a phoneme inventory --- but that inventory is derived, not
assumed.
\end{remark}

\subsection{Linguistics Observable Family}\label{sec:ling-obs-family}

\begin{definition}[\status{D}]\label{def:ling-obs-family}
\textup{(Linguistics Observable Family.)}
A \emph{linguistics observable family} on $\mathcal{R}_{\mathrm{Ling}}$
is a triple
\[
  \mathcal{O}_{\mathrm{Ling}} \;=\;
  \bigl(K_{\mathrm{Ling}},\;
        \mathcal{C}_{\mathrm{Ling}},\;
        \Gamma_{\mathrm{Ling}}\bigr)
\]
where:
\begin{enumerate}[label=\textup{(\arabic*)}]
  \item $K_{\mathrm{Ling}}$ is the \emph{contrast kernel}: the
    quotient-visible distinguishability data induced on
    $\mathcal{C}_{\mathrm{obs}}$ by the admissible contrast discipline
    $\mathcal{P}$; it is the analogue of $\mathrm{Spec}(L_{\mathrm{Spec}})$
    in the SPEC branch;
  \item $\mathcal{C}_{\mathrm{Ling}}$ is the \emph{compositional ledger}:
    certified records of one-step assembly operations and their
    transport costs, analogous to the transition ledger of the SPEC
    branch; and
  \item $\Gamma_{\mathrm{Ling}}$ is the \emph{grammar-descent object}:
    a branch-visible certified rule object governing when an assembled
    sequence of contrast-class representatives is admissible within the
    declared compositional window.
\end{enumerate}
\end{definition}

\begin{definition}[\status{D}]\label{def:ling-contrast-kernel}
\textup{(Contrast Kernel.)}
The \emph{contrast kernel} $K_{\mathrm{Ling}}$ is the collection of
ordered pairs $(X, X')$ of signal-objects together with the certified
set of admissible contrast probes that distinguish them, modulo the
observational quotient imposed by Book~I's equivalence discipline.
Two signal-objects belong to the same \emph{contrast equivalence class}
if and only if no admissible probe in $\mathcal{P}$ distinguishes them.
\end{definition}

\begin{definition}[\status{D}]\label{def:ling-contrast-equivalence}
\textup{(Contrast Equivalence Class.)}
A \emph{contrast equivalence class} $[X]_{K}$ is the quotient class of
a signal-object $X$ under the contrast kernel $K_{\mathrm{Ling}}$.
The collection of all contrast equivalence classes on the declared
regime is the \emph{certified contrast inventory} --- the canonical
analogue of a phoneme or morpheme inventory, derived rather than
assumed.
\end{definition}

\begin{remark}[\status{R}]\label{rem:ling-phoneme-remark}
The certified contrast inventory is not assumed to coincide with any
pre-existing phonological category system.  If, under a given regime
and admissible contrast discipline, the contrast equivalence classes
stabilize to a finite partition matching a known phoneme inventory,
that match is a derived theorem, not a branch assumption.
\end{remark}

\begin{definition}[\status{D}]\label{def:ling-contrast-margin}
\textup{(Contrast Margin.)}
The \emph{contrast margin}
\[
  \Delta_{\mathrm{Ling}}(X, X')
\]
is a branch-specific nonnegative quantity measuring the observable
separation between two contrast equivalence classes $[X]_{K}$ and
$[X']_{K}$ within the declared regime.  It is positive if and only if
$[X]_{K} \neq [X']_{K}$.  This is the linguistic analogue of the
spectral margin of the SPEC branch and the mass gap of the YM branch.
\end{definition}

\begin{definition}[\status{D}]\label{def:ling-compositional-ledger}
\textup{(Compositional Ledger.)}
The \emph{compositional ledger} $\mathcal{C}_{\mathrm{Ling}}$ is a
certified bookkeeping object assigning to each admissible one-step
assembly operation --- combining two contrast-class representatives
into a larger signal-object --- the induced change in branch-visible
contrast signature, together with any declared transport cost or visible
defect contribution.
\end{definition}

\begin{remark}[\status{R}]\label{rem:ling-compositionality-remark}
The compositional ledger is a probe-ledger for assembly operations.
It does not assume that the result of an assembly is a sentence, a
phrase, or any other pre-theoretic linguistic unit.  Whether such units
emerge as stable quotient-visible objects is a derived question, not a
branch primitive.
\end{remark}

\begin{definition}[\status{D}]\label{def:ling-grammar-descent}
\textup{(Grammar-Descent Object.)}
A \emph{grammar-descent object} $\Gamma_{\mathrm{Ling}}$ is a
branch-visible certified rule object that:
\begin{enumerate}[label=\textup{(\arabic*)}]
  \item is sourced through $\mathbf{U}$;
  \item governs which assembly sequences are admissible within the
    declared compositional window; and
  \item produces only quotient-visible distinctions among assembled
    objects.
\end{enumerate}
$\Gamma_{\mathrm{Ling}}$ is not a context-free grammar, a minimalist
feature system, or any other pre-given formalism.  It is a derived
certified object, admitted only when the branch-legitimacy conditions of
Book~V are satisfied.
\end{definition}

\subsection{Endpoint Datum}\label{sec:ling-endpoint}

\begin{definition}[\status{D}]\label{def:ling-endpoint}
\textup{(Linguistics Endpoint Datum.)}
The \emph{linguistics endpoint datum} is the branch-visible package
\[
  E_{\mathrm{Ling}} \;=\;
  \bigl(\mathcal{C}_{\mathrm{obs}},\;
        \mathcal{O}_{\mathrm{Ling}},\;
        \Delta_{\mathrm{Ling}},\;
        \mathcal{W}_{\mathrm{comp}}\bigr)
\]
consisting of the certified observable class, the linguistics
observable family, the contrast margin, and the declared compositional
window.
\end{definition}

\begin{definition}[\status{D}]\label{def:lawful-ling-descendant}
\textup{(Lawful Linguistics Descendant Claim.)}
A \emph{lawful linguistics descendant claim} is a claim that the
endpoint datum $E_{\mathrm{Ling}}$ is sourced through $\mathbf{U}$ and
supports a certified contrast, compositionality, or grammar-descent
statement on the declared regime.
\end{definition}

\subsection{Current Branch Posture}\label{sec:ling-posture}

\begin{proposition}[\status{C}]\label{prop:ling-branch-candidate}
\textup{(Linguistics Branch Candidate.)}
\textup{[Conditional on the linguistics target regime, observable
family, and endpoint datum all being sourced through $\mathbf{U}$,
branch-visible, and admitting a Book~V legitimacy witness.]}
The linguistics program is a lawful branch candidate.
\end{proposition}

\begin{proof}
Immediate from the branch-legitimacy doctrine (Book~V) applied to the
declared linguistics endpoint datum.
\end{proof}

\begin{remark}[\status{R}]\label{rem:ling-posture-note}
Prop.~\ref{prop:ling-branch-candidate} constitutes the branch
architecturally.  It does not yet certify any contrast theorem,
compositional theorem, or grammar-descent theorem.
\end{remark}

%% ---------------------------------------------------------------
\section{Admissible Contrast Probes and Contrast Signatures}
\label{sec:ling-probes}
%% ---------------------------------------------------------------

\subsection{Admissible Contrast Probes}\label{sec:ling-admissible-probes}

\begin{definition}[\status{D}]\label{def:ling-admissible-probe}
\textup{(Admissible Contrast Probe.)}
An \emph{admissible contrast probe} is a declared admissible process
$P \in \mathcal{P}$ whose effect on the certified observable class is
branch-visible and accountable within the spine's admissible-test
discipline.  In the linguistic setting, paradigm cases include minimal
contrast tests, substitution tests, and admissible distributional
comparison tests --- but none of these is assumed as a primitive; each
must be declared and certified within the NFC framework.
\end{definition}

\begin{standingrule}[\status{D}]\label{sr:ling-no-hidden-probe}
No admissible contrast probe may introduce hidden feature labels,
undeclared phonological representations, or non-certified side
channels.  Any probe violating this rule is not admissible for
linguistics branch claims.
\end{standingrule}

\begin{definition}[\status{D}]\label{def:ling-probe-episode}
\textup{(Contrast Probe Episode.)}
A \emph{contrast probe episode} on a pair of signal-objects
$(X, X') \in \mathcal{C}_{\mathrm{obs}} \times \mathcal{C}_{\mathrm{obs}}$
is an ordered application of an admissible contrast probe together with
the certified before/after distinction record visible in the declared
regime.
\end{definition}

\begin{remark}[\status{R}]\label{rem:ling-probe-episode-note}
A contrast probe episode is the linguistics analogue of a controlled
admissible comparison, not yet a phonological judgment event in the
ordinary external sense.
\end{remark}

\subsection{Contrast Signatures}\label{sec:ling-contrast-signatures}

\begin{definition}[\status{D}]\label{def:ling-contrast-signature}
\textup{(Contrast Signature.)}
A \emph{contrast signature} is a branch-visible assignment
\[
  \mathsf{Resp}_{L}(X; P)
\]
from a signal-object $X$ and admissible probe $P$ to a quotient-visible
detection outcome, together with its transport and defect bookkeeping.
The contrast signature is the linguistic analogue of the spectroscopic
response signature.
\end{definition}

\begin{standingrule}[\status{D}]\label{sr:ling-signature-invariance}
A contrast signature has theorem-bearing force only to the extent that
it is invariant under presentation-level differences already eliminated
by the observational quotient.  Two signal-objects that are
observationally equivalent cannot bear different contrast signatures.
\end{standingrule}

\begin{definition}[\status{D}]\label{def:ling-contrast-equiv-signature}
\textup{(Contrast-Signature Equivalence.)}
Two contrast signatures are \emph{contrast-equivalent} if they
determine the same quotient-visible detection outcome class under the
declared comparison discipline of $\mathcal{R}_{\mathrm{Ling}}$.
\end{definition}

\subsection{One-Step Compositional Ledger}\label{sec:ling-one-step}

\begin{definition}[\status{D}]\label{def:ling-one-step-comp}
\textup{(One-Step Compositional Ledger Entry.)}
Given an admissible assembly operation combining $X$ and $X'$ into an
assembled signal-object $X \cdot X'$ (within the declared compositional
window), the \emph{one-step compositional ledger entry}
\[
  \Delta_{\mathrm{comp}}\!\bigl(X,\, X'\bigr)
\]
records the certified contrast-class of $X \cdot X'$, together with
any declared transport cost or visible defect contribution inherited
from $X$ and $X'$.
\end{definition}

\begin{remark}[\status{R}]\label{rem:ling-assembly-note}
The assembly operation $X \cdot X'$ is not concatenation in the
ordinary sense.  It is any admissible process (declared within the
branch regime) that produces a new signal-object from two existing
ones, with all transport bookkeeping declared.  Whether a familiar
notion of string concatenation is lawfully licensed here is a
downstream derived question.
\end{remark}

%% ---------------------------------------------------------------
\section{Core Theorem Shells and Named Frontiers}
\label{sec:ling-theorems}
%% ---------------------------------------------------------------

The following theorem shells articulate the main claims of the
linguistics branch.  Each carries the status it would bear
\emph{if} its stated open obligations were discharged.  The open
obligations are named explicitly as branch frontiers.

\subsection{Contrast Realization}\label{sec:ling-CR}

\begin{openobligation}[\status{O}]\label{ob:ling-CR}
\textup{(O-LING.CR: Contrast Realization.)}
Show that the contrast kernel $K_{\mathrm{Ling}}$ is sourced through
$\mathbf{U}$ and constitutes a lawful branch observable within the
declared regime, with all probe declarations inherited from the
Book~I admissible-test discipline.
\end{openobligation}

\begin{theorem}[\status{C}]\label{thm:ling-contrast-realization}
\textup{(Contrast Realization Theorem.)}
\textup{[Conditional on O-LING.CR: Contrast Realization.]}
The contrast kernel $K_{\mathrm{Ling}}$ is a lawful branch observable
descended from $\mathbf{U}$.  The certified contrast equivalence
classes $\{[X]_K\}$ are quotient-visible and admissibly testable.
\end{theorem}

\begin{proof}
By the branch-legitimacy doctrine (Book~V) applied to the declared
contrast kernel, together with the Book~I observational quotient
discipline.  The full proof is conditional on O-LING.CR.
\end{proof}

\subsection{Compositional Ledger Theorem}\label{sec:ling-CL}

\begin{openobligation}[\status{O}]\label{ob:ling-CL}
\textup{(O-LING.CL: Compositional Ledger.)}
Show that the compositional ledger $\mathcal{C}_{\mathrm{Ling}}$ is
additive, sourced through the Book~III transport accounting, and that
one-step assembly operations preserve transport bookkeeping without
hidden defect.
\end{openobligation}

\begin{theorem}[\status{C}]\label{thm:ling-comp-ledger}
\textup{(Compositional Ledger Theorem.)}
\textup{[Conditional on O-LING.CL: Compositional Ledger.]}
The compositional ledger $\mathcal{C}_{\mathrm{Ling}}$ is additive
over declared assembly steps: transport costs distribute across
components without hidden channel.
\end{theorem}

\begin{proof}
By the Book~III unified coercive-transport inequality applied to the
assembly operation, together with the no-hidden-channel standing rule
(Sr.~\ref{sr:ling-no-hidden-probe}).  Full proof is conditional on
O-LING.CL.
\end{proof}

\subsection{Contrast Margin Theorem}\label{sec:ling-CM}

\begin{openobligation}[\status{O}]\label{ob:ling-CM}
\textup{(O-LING.CM: Contrast Margin.)}
Show that the contrast margin $\Delta_{\mathrm{Ling}}$ is
branch-computably positive across distinct contrast equivalence classes
within the declared regime, and that no equivalence class collapse
occurs without a declared admissible merger.
\end{openobligation}

\begin{theorem}[\status{C}]\label{thm:ling-contrast-margin}
\textup{(Contrast Margin Theorem.)}
\textup{[Conditional on O-LING.CM: Contrast Margin.]}
For any two distinct contrast equivalence classes $[X]_K \neq [X']_K$
in the certified contrast inventory, $\Delta_{\mathrm{Ling}}(X,X') > 0$.
The contrast inventory is discretely separated on the declared regime.
\end{theorem}

\begin{proof}
By the observational quotient discipline (Book~I) and the declared
admissible-test finiteness condition on $\mathcal{R}_{\mathrm{Ling}}$.
The positivity of the contrast margin reflects the separating power of
the admissible contrast discipline.  Full proof is conditional on
O-LING.CM.
\end{proof}

\begin{remark}[\status{R}]\label{rem:ling-gap-analogy}
The contrast margin theorem is the linguistics branch analogue of the
mass gap theorem (YM branch) and the spectral margin theorem (SPEC
branch).  In each case the key structural fact is a positive separation
between inequivalent observable classes.  Here the separation is
between phonologically (or otherwise observationally) distinct
contrast equivalence classes.
\end{remark}

\subsection{No-Hidden-Channel Theorem}\label{sec:ling-NC}

\begin{openobligation}[\status{O}]\label{ob:ling-NC}
\textup{(O-LING.NC: No Hidden Channel.)}
Show that no compositional assembly pathway exists outside the
declared one-step compositional ledger: any discrepancy between two
assembled objects sharing the same declared component contrast classes
would require an undeclared assembly channel, ruled out by the
exhaustive source-image condition.
\end{openobligation}

\begin{theorem}[\status{C}]\label{thm:ling-no-hidden-channel}
\textup{(No-Hidden-Channel Theorem.)}
\textup{[Conditional on O-LING.NC: No Hidden Channel.]}
Two assembled signal-objects sharing the same declared component
contrast classes and the same admissible assembly operation have the
same contrast signature.  No undeclared compositional pathway can
produce a branch-visible distinction invisible to the declared ledger.
\end{theorem}

\begin{proof}
Any discrepancy would require an undeclared dissipative or
combinatorial channel.  By the exhaustive source-image condition
(\texttt{def:exhaustive-source-image}, Book~IV) applied to
$\mathbf{U}$ and the no-hidden-symbol standing rule
(Sr.~\ref{sr:ling-no-hidden-symbol}), no such channel is admitted.
Full proof is conditional on O-LING.NC.
\end{proof}

\subsection{Grammar Descent Theorem}\label{sec:ling-GD}

\begin{openobligation}[\status{O}]\label{ob:ling-GD}
\textup{(O-LING.GD: Grammar Descent.)}
Show that the grammar-descent object $\Gamma_{\mathrm{Ling}}$ is
sourced through $\mathbf{U}$, passes the five-condition Book~V
branch-legitimacy test, and produces only quotient-visible distinctions
among assembled objects within the declared compositional window.
\end{openobligation}

\begin{theorem}[\status{C}]\label{thm:ling-grammar-descent}
\textup{(Grammar Descent Theorem.)}
\textup{[Conditional on O-LING.GD: Grammar Descent.]}
The grammar-descent object $\Gamma_{\mathrm{Ling}}$ is a lawful
branch-visible certified rule object.  Its admissible assembly
sequences are exactly those whose contrast-class outcomes are
branch-visible within $\mathcal{R}_{\mathrm{Ling}}$.
\end{theorem}

\begin{proof}
By the Book~V branch-legitimacy doctrine applied to
$\Gamma_{\mathrm{Ling}}$, together with the no-hidden-symbol and
no-smuggling standing rules.  Full proof is conditional on O-LING.GD.
\end{proof}

\subsection{Recursion and Infinite Generativity}\label{sec:ling-REC}

\begin{openobligation}[\status{O}]\label{ob:ling-REC}
\textup{(O-LING.REC: Recursion.)}
Show, or certify the impossibility of showing within the NFC framework,
that the compositional ledger $\mathcal{C}_{\mathrm{Ling}}$ supports
iterated self-embedding without bound --- the formal analogue of
linguistic recursion.  This obligation is expected to require a
Book~VI bridge theorem licensing the infinite-generativity notion as
a certified asymptotic encoding of finite iterated assembly, and may
be permanently open at the finite-observable level.
\end{openobligation}

\begin{remark}[\status{R}]\label{rem:ling-rec-posture}
\textbf{Posture on recursion.}
Recursion in the ordinary linguistic sense refers to the unbounded
self-embedding of syntactic constituents.  Within the NFC framework
this cannot be a finite observable fact --- any finite probe episode
sees only finitely many nesting levels.  If recursion appears at all,
it must be licensed by a Book~VI bridge theorem as the asymptotic limit
of a certified sequence of finite assembly operations.  Whether such a
bridge theorem can be proved is a genuine open question, not a defect
of the branch architecture.
\end{remark}

\subsection{Semantic Grounding}\label{sec:ling-SEM}

\begin{openobligation}[\status{O}]\label{ob:ling-SEM}
\textup{(O-LING.SEM: Semantic Grounding.)}
Show that a semantic assignment --- a quotient-visible branch object
playing the role of meaning for contrast-class sequences --- can be
sourced through $\mathbf{U}$, without importing a background semantic
domain, truth-value space, or model-theoretic universe.  Semantic
grounding is expected to require the full grammar-descent result plus
a certified probe-response structure on the assembled objects.
\end{openobligation}

\begin{remark}[\status{R}]\label{rem:ling-sem-posture}
\textbf{Posture on semantics.}
Semantic grounding is the deepest downstream obligation in this branch.
No claim of meaning is made here.  The no-semantic-primitive standing
rule (Sr.~\ref{sr:ling-no-semantic-primitive}) prohibits importing a
semantic domain.  The task of showing that a quotient-visible semantic
assignment exists and is sourced through $\mathbf{U}$ is a named open
obligation, not a currently certified result.
\end{remark}

\subsection{Linguistics Endpoint}\label{sec:ling-END}

\begin{openobligation}[\status{O}]\label{ob:ling-END}
\textup{(O-LING.END: Linguistics Endpoint.)}
Show that the endpoint datum $E_{\mathrm{Ling}}$ is fully certified:
contrast realization, compositional ledger, contrast margin, no-hidden
channel, and grammar descent all discharged.  The endpoint is reached
when the linguistics branch satisfies the shared-endgame schema of
Book~VII.
\end{openobligation}

\begin{theorem}[\status{C}]\label{thm:ling-endpoint}
\textup{(Linguistics Endpoint Theorem.)}
\textup{[Conditional on O-LING.CR, O-LING.CL, O-LING.CM, O-LING.NC,
O-LING.GD.]}
The endpoint datum
\[
  E_{\mathrm{Ling}} \;=\;
  \bigl(\mathcal{C}_{\mathrm{obs}},\;
        \mathcal{O}_{\mathrm{Ling}},\;
        \Delta_{\mathrm{Ling}},\;
        \mathcal{W}_{\mathrm{comp}}\bigr)
\]
is a certified lawful linguistics endpoint: certified contrast
structure, additive compositional ledger, positive contrast margin,
no-hidden-channel control, and grammar-descent sourced through
$\mathbf{U}$.  The contrast-founded compositionality endpoint is
reached under the stated conditions.
\end{theorem}

\begin{proof}
Combine Thms.~\ref{thm:ling-contrast-realization},
\ref{thm:ling-comp-ledger}, \ref{thm:ling-contrast-margin},
\ref{thm:ling-no-hidden-channel}, and
\ref{thm:ling-grammar-descent}.  Apply the shared-endgame schema
(Book~VII, Prop.~6.11): part~(i) is the source-controlled realized
object given by the contrast kernel descended from $\mathbf{U}$;
part~(ii) is the compositional ledger theorem; part~(iii) is the
no-hidden-channel theorem together with the contrast margin.  Full
proof is conditional on O-LING.CR through O-LING.GD.
\end{proof}

%% ---------------------------------------------------------------
\section{Closure Ledger}\label{sec:ling-closure}
%% ---------------------------------------------------------------

\subsection{Summary of Named Open Obligations}

The following table records all named open obligations in the
linguistics branch as of the current canonical state.

\medskip
\noindent
\begin{tabular}{|l|l|p{8.5cm}|}
\hline
\textbf{Label} & \textbf{Status} & \textbf{Description} \\
\hline
\texttt{ob:ling-CR} & [O] & Contrast Realization: $K_{\mathrm{Ling}}$ sourced through $\mathbf{U}$ \\
\texttt{ob:ling-CL} & [O] & Compositional Ledger: additivity and transport accounting \\
\texttt{ob:ling-CM} & [O] & Contrast Margin: positive separation of distinct classes \\
\texttt{ob:ling-NC} & [O] & No Hidden Channel: no undeclared compositional pathway \\
\texttt{ob:ling-GD} & [O] & Grammar Descent: $\Gamma_{\mathrm{Ling}}$ passes Book~V test \\
\texttt{ob:ling-REC} & [O] & Recursion: infinite generativity via Book~VI bridge (may remain open) \\
\texttt{ob:ling-SEM} & [O] & Semantic Grounding: meaning as quotient-visible branch object \\
\texttt{ob:ling-END} & [O] & Endpoint: full certification of $E_{\mathrm{Ling}}$ \\
\hline
\end{tabular}

\medskip

\begin{remark}[\status{R}]\label{rem:ling-discharge-order}
\textbf{Discharge order.}
The natural discharge order is: CR $\to$ CL $\to$ CM $\to$ NC $\to$
GD $\to$ END.  Recursion (REC) and Semantic Grounding (SEM) are
downstream of END and may remain permanently open.  Crucially, END can
be reached without discharging REC or SEM: the
contrast-founded-compositionality endpoint is the minimal honest claim
of this branch.
\end{remark}

\subsection{Two-Status Reading}

\begin{remark}[\status{R}]\label{rem:ling-two-status}
\textbf{Two-status reading.}
The linguistics branch admits two natural closure readings:
\begin{enumerate}[label=\textup{(\alph*)}]
  \item \emph{Contrast-compositionality closure} (obligations CR--END
    discharged, REC and SEM open): certified contrast inventory,
    additive compositional ledger, positive contrast margin, and a
    grammar-descent object --- the minimal honest linguistic endpoint
    within the NFC framework.
  \item \emph{Semantic closure} (obligations CR--SEM discharged): adds
    a quotient-visible meaning assignment, licensed as a probe-response
    structure on assembled objects.  Requires the full grammar-descent
    result and a certified probe-response object at the compositional
    level.
\end{enumerate}
Posture~(a) is the correct near-term target.  Posture~(b) is the
genuine long-term ambition.
\end{remark}

\subsection{Relation to Existing Branches}

\begin{remark}[\status{R}]\label{rem:ling-branch-relations}
\textbf{Relation to existing branches.}
The linguistics branch imports structural templates from the SPEC and
NS branches but does not depend on any of their open obligations being
discharged.  In particular:
\begin{enumerate}[label=\textup{(\arabic*)}]
  \item The contrast-probe structure is formally parallel to the
    probe-response structure of the SPEC branch (\texttt{def:probe-response-object},
    \texttt{def:transition-ledger}), but operates on signal-object
    pairs rather than physical resonance data.  No SPEC branch
    obligation is a prerequisite.
  \item The compositional propagation structure is formally parallel to
    the NS continuation criterion, but governed by admissible assembly
    rather than differential-equation dynamics.  No NS branch
    obligation is a prerequisite.
  \item The contrast margin is the linguistic analogue of the YM mass
    gap (positive spectral margin) and the SPEC spectral margin.  No YM
    or SPEC margin theorem is a prerequisite; the linguistic instance is
    derived independently from the admissible contrast discipline.
  \item The SM matter branch and the RH branch have no direct bearing
    on the linguistics branch at the contrast-compositionality level.
    Cross-branch obligations may emerge if semantics requires matter
    content, but this is not anticipated at the current architectural
    stage.
\end{enumerate}
\end{remark}

%% ---------------------------------------------------------------
\section{Obligation Discharges}\label{sec:ling-discharges}
%% ---------------------------------------------------------------

This section discharges obligations O-LING.CR through O-LING.END by
appeal to existing certified spine results.  Each discharge replaces
the corresponding open obligation with a conditional theorem whose
hypotheses are stated explicitly.  O-LING.REC and O-LING.SEM remain
open obligations: recursion is genuinely open (expected permanently
open at the finite-observable level) and semantic grounding has no
discharge path available in the current canon.

\subsection{O-LING.CR: Contrast Realization}\label{sec:ling-CR-discharge}

\begin{theorem}[\status{C}]\label{thm:ling-CR-discharged}
\textup{(Contrast Realization --- O-LING.CR Discharged.)}
\textup{[Conditional on: (1) a declared linguistics target regime
$\mathcal{R}_{\mathrm{Ling}}$ with admissible contrast probe family
$\mathcal{P} \subseteq \mathcal{T}_n$ (Book~I admissible tests) and
certified observable class $\mathcal{C}_{\mathrm{obs}}$ descended from
$\mathbf{U}$ (Book~IV, \texttt{thm:universal-completion}); and (2)
Book~II transfer structure (\texttt{def:transfer-system}) on
$\mathcal{C}_{\mathrm{obs}}$.]}
There exists a certified contrast kernel
\[
  K_{\mathrm{Ling}} : (X, X') \;\mapsto\;
  \bigl([X]_{\sim_n},\, [X']_{\sim_n},\,
  \{P \in \mathcal{P} : P \text{ distinguishes } X \text{ from } X'\}\bigr)
\]
where $[\cdot]_{\sim_n}$ is the Book~I observable equivalence class.
The contrast equivalence classes $\{[X]_K\}$ are quotient-visible,
admissibly testable, and transport-accountable under Book~III.
\end{theorem}

\begin{proof}
The observable equivalence relation $\sim_n$ (Book~I,
\texttt{def:obs-equiv}) is already defined on all admissible-test
outcomes.  Restricting to the contrast probe family
$\mathcal{P} \subseteq \mathcal{T}_n$ gives quotient-visible contrast
equivalence classes by the no-smuggling standing rule
(Sr.~\ref{sr:ling-no-smuggling}).  The contrast kernel records exactly
the probes that witness distinctions, together with the equivalence
classes of the two signal-objects.  Transport accounting is provided
by Book~III persistence and transfer structure.  The
boundary-capacity bookkeeping from Book~II bounds the information
capacity of any contrast boundary within $\mathcal{R}_{\mathrm{Ling}}$.
No imported symbol system or phonological primitive is required.
\qed
\end{proof}

\begin{remark}[\status{R}]\label{rem:ling-CR-scope}
The contrast kernel here is the structural minimum: it records
observable equivalence classes and the probes that separate them.
The richer structure --- compositional ledger, grammar descent,
contrast margin --- builds on this foundation in subsequent discharges.
\end{remark}

\subsection{O-LING.CL: Compositional Ledger}\label{sec:ling-CL-discharge}

\begin{theorem}[\status{C}]\label{thm:ling-CL-discharged}
\textup{(Compositional Ledger --- O-LING.CL Discharged.)}
\textup{[Conditional on: O-LING.CR discharged
(Thm.~\ref{thm:ling-CR-discharged}); Book~III unified
coercive-transport inequality (\texttt{thm:unified-coercive}, [U])
and ledger additivity (\texttt{lem:transport-equiv-null}).]}
The compositional ledger $\mathcal{C}_{\mathrm{Ling}}$ is additive over
declared assembly steps.  Concretely, for a one-step assembly of
$X$ and $X'$ into $X \cdot X'$, the ledger entry is
\[
  \Delta_{\mathrm{comp}}\!\bigl(X,\, X'\bigr)
  \;:=\; [X \cdot X']_{\sim_n}
\]
with transport cost bounded by the Book~III coercive-transport
inequality applied to the assembly operation.  Multi-step assembly
costs distribute additively without hidden channel.
\end{theorem}

\begin{proof}
Both $[X]_{\sim_n}$ and $[X']_{\sim_n}$ are quotient-visible by
O-LING.CR.  The assembly operation $X \cdot X'$ is a declared
admissible process; its observable equivalence class is
$[X \cdot X']_{\sim_n}$.  By the Book~III unified coercive-transport
inequality, the transport cost of the assembly is bounded by the sum
of the constituent costs plus any declared defect, and by
\texttt{lem:transport-equiv-null} the ledger entry is invariant under
admissible relabelling of the components.  Additivity over multi-step
assembly follows by induction, with each step contributing its
transport cost to the running total.  The no-hidden-probe standing
rule (Sr.~\ref{sr:ling-no-hidden-probe}) ensures no side channels
enter.
\qed
\end{proof}

\subsection{O-LING.CM: Contrast Margin}\label{sec:ling-CM-discharge}

\begin{theorem}[\status{C}]\label{thm:ling-CM-discharged}
\textup{(Contrast Margin --- O-LING.CM Discharged.)}
\textup{[Conditional on: O-LING.CR discharged
(Thm.~\ref{thm:ling-CR-discharged}); the admissible contrast probe
family $\mathcal{P}$ being non-trivially separating on the declared
regime, i.e., for any two distinct equivalence classes $[X]_K \neq
[X']_K$ there exists at least one $P \in \mathcal{P}$ distinguishing
them.]}
For any two distinct contrast equivalence classes $[X]_K \neq [X']_K$
in the certified contrast inventory, $\Delta_{\mathrm{Ling}}(X,X') > 0$.
The certified contrast inventory is discretely separated on the
declared regime.
\end{theorem}

\begin{proof}
By definition, $[X]_K \neq [X']_K$ if and only if some probe
$P \in \mathcal{P}$ distinguishes $X$ from $X'$ within the declared
regime.  The contrast margin $\Delta_{\mathrm{Ling}}(X, X')$ is defined
as the quotient-visible measure of this separation; it is positive
whenever at least one separating probe exists
(Def.~\ref{def:ling-contrast-margin}).  The non-trivially-separating
condition precisely requires this: no two distinct equivalence classes
are probe-indistinguishable within $\mathcal{R}_{\mathrm{Ling}}$.
Therefore $\Delta_{\mathrm{Ling}}(X,X') > 0$ for all distinct pairs.
\qed
\end{proof}

\begin{remark}[\status{R}]\label{rem:ling-CM-note}
The contrast margin theorem is essentially definitional once the
contrast probe family is declared to be non-trivially separating.  The
substantive work is in establishing which probe family satisfies this
condition for a given physical or formal language target.  The theorem
records the structural consequence cleanly: a declared separating
probe family implies a positive contrast margin.
\end{remark}

\subsection{O-LING.NC: No Hidden Channel}\label{sec:ling-NC-discharge}

\begin{theorem}[\status{C}]\label{thm:ling-NC-discharged}
\textup{(No Hidden Compositional Channel --- O-LING.NC Discharged.)}
\textup{[Conditional on: O-LING.CL discharged
(Thm.~\ref{thm:ling-CL-discharged}); two assembled signal-objects
sharing the same declared component contrast classes and the same
admissible assembly operation; and the universal completion theorem
(\texttt{thm:universal-completion}, Book~IV, [U]) applied via the
exhaustive source-image condition
(\texttt{def:exhaustive-source-image}, Book~IV).]}
The two assembled signal-objects have the same contrast signature.
No undeclared compositional pathway can produce a branch-visible
distinction invisible to the declared ledger.
\end{theorem}

\begin{proof}
Let $A = X \cdot X'$ and $B = Y \cdot Y'$ where $[X]_K = [Y]_K$,
$[X']_K = [Y']_K$, and the same admissible assembly operation is
applied.  By O-LING.CL, the compositional ledger entry for each is
determined by the equivalence classes $[X]_K$, $[X']_K$ and the
transport cost, both of which are shared.  Suppose for contradiction
that $[A]_K \neq [B]_K$.  Then some probe $P \in \mathcal{P}$
distinguishes $A$ from $B$.  This probe witnesses a difference in
their observable equivalence classes that is not attributable to
any declared component-level or assembly-level difference.  Such a
difference would require an undeclared compositional channel --- a
source of observable variation not appearing in the ledger.  By the
exhaustive source-image condition
(\texttt{def:exhaustive-source-image}, Book~IV) applied to
$\mathbf{U}$, no such undeclared channel exists.  Contradiction.
Hence $[A]_K = [B]_K$ and the contrast signatures agree.
\qed
\end{proof}

\subsection{O-LING.GD: Grammar Descent}\label{sec:ling-GD-discharge}

\begin{theorem}[\status{C}]\label{thm:ling-GD-discharged}
\textup{(Grammar Descent --- O-LING.GD Discharged.)}
\textup{[Conditional on: O-LING.CR (Thm.~\ref{thm:ling-CR-discharged})
and O-LING.NC (Thm.~\ref{thm:ling-NC-discharged}) discharged; the
declared admissible contrast discipline $\mathcal{P}$ being closed
under the declared assembly operation in the sense that probe outcomes
on assembled objects are determined by probe outcomes on components;
and Book~V branch-legitimacy doctrine (\texttt{thm:UBLT}, [U]).]}
There exists a certified grammar-descent object $\Gamma_{\mathrm{Ling}}$
sourced through $\mathbf{U}$, identified as the restriction of the
Book~III coercive-transport rule to the contrast-visible sector of
$\mathcal{C}_{\mathrm{obs}}$ on $\mathcal{R}_{\mathrm{Ling}}$:
\begin{enumerate}[label=\textup{(\arabic*)}]
  \item $\Gamma_{\mathrm{Ling}}$ governs admissible assembly sequences
    within the declared compositional window;
  \item its outputs are quotient-visible within $\mathcal{R}_{\mathrm{Ling}}$;
  \item it passes the Book~V five-condition branch-legitimacy test;
  \item the admissible assembly sequences it licenses are exactly
    those whose contrast-class outcomes are branch-visible.
\end{enumerate}
\end{theorem}

\begin{proof}
By O-LING.CR, the observable equivalence classes on
$\mathcal{C}_{\mathrm{obs}}$ are certified.  The Book~III
coercive-transport governing theorem (\texttt{thm:governing}) constructs
a canonical rule on every certified transport-system descendant of
$\mathbf{U}$.  Restricting to $\mathcal{C}_{\mathrm{obs}}$ on
$\mathcal{R}_{\mathrm{Ling}}$ gives $\Gamma_{\mathrm{Ling}}$.  By
O-LING.NC, the assembly outputs of $\Gamma_{\mathrm{Ling}}$ are
uniquely determined by their component contrast classes: no hidden
channel breaks the linkage.  The admissible assembly sequences
licensed by $\Gamma_{\mathrm{Ling}}$ are therefore exactly those
producing branch-visible contrast-class outcomes.  The Book~V
five-condition test is satisfied:
\begin{enumerate}[label=\textup{(\roman*)}]
  \item explicit branch candidate ($\Gamma_{\mathrm{Ling}}$ declared);
  \item observable family visibility (outputs are quotient-visible);
  \item source-descended comparison (inherits from $\mathbf{U}$ via
    Book~III);
  \item no hidden supplement (Sr.~\ref{sr:ling-no-hidden-symbol});
  \item no path multiplication (O-LING.NC ensures uniqueness of
    contrast-class output per assembly operation).
\end{enumerate}
\qed
\end{proof}

\subsection{O-LING.END: Linguistics Endpoint}\label{sec:ling-END-discharge}

\begin{theorem}[\status{C}]\label{thm:ling-END-discharged}
\textup{(Linguistics Endpoint --- O-LING.END Discharged.)}
\textup{[Conditional on O-LING.CR through O-LING.GD all discharged
(Thms.~\ref{thm:ling-CR-discharged}--\ref{thm:ling-GD-discharged}); the
BR source theorem (\texttt{thm:BR-source}, Book~V, [C]); and the ICDC
schema (\texttt{thm:ICDC}, Book~V, [C]).]}
The endpoint datum
\[
  E_{\mathrm{Ling}} \;=\;
  \bigl(\mathcal{C}_{\mathrm{obs}},\;
        \mathcal{O}_{\mathrm{Ling}},\;
        \Delta_{\mathrm{Ling}},\;
        \mathcal{W}_{\mathrm{comp}}\bigr)
\]
is a certified lawful linguistics descendant closure claim on the declared
contrast-compositionality regime:
\begin{enumerate}[label=\textup{(\roman*)}]
  \item branch-visible contrast equivalence classes are certified
    (O-LING.CR);
  \item compositional ledger is additive and transport-accountable
    (O-LING.CL);
  \item contrast inventory is discretely separated (O-LING.CM);
  \item no undeclared compositional pathway exists (O-LING.NC);
  \item grammar-descent object passes Book~V legitimacy test
    (O-LING.GD);
  \item the entire package is sourced through $\mathbf{U}$ via the BR
    coercive-transport route and is ICDC-coherent.
\end{enumerate}
\end{theorem}

\begin{proof}
Items~(i)--(v) follow directly from the named discharged obligations.
Item~(vi): by O-LING.GD, $\Gamma_{\mathrm{Ling}}$ is the restriction
of the Book~III coercive operator to the contrast-visible sector, which
the BR source theorem (\texttt{thm:BR-source}) certifies as a lawful
specialization of the upstream coercive-transport form on $\mathbf{U}$.
The ICDC schema applies with no-hidden-channel control as the one-step
inheritance datum: the gap-subcritical ICDC reading ($\Theta_n = 1$)
gives the contrast-compositionality-local statement that certified
contrast structure persists under admissible assembly.  Apply the
Book~VII shared-endgame schema (Prop.~6.11): part~(i) is the
source-controlled contrast kernel descended from $\mathbf{U}$;
part~(ii) is the compositional ledger theorem; part~(iii) is the
no-hidden-channel theorem together with the contrast margin.  The
endpoint datum is therefore fully assembled, scoped, and replayable.
\qed
\end{proof}

\begin{remark}[\status{R}]\label{rem:ling-contrast-comp-complete}
\textbf{Contrast-compositionality endpoint: discharge status.}
Under the hypotheses of
Thms.~\ref{thm:ling-CR-discharged}--\ref{thm:ling-END-discharged},
the contrast-compositionality program reaches its endpoint.  The
branch carries: a certified contrast kernel; additive compositional
ledger; positive contrast margin; no-hidden-channel control; and a
grammar-descent object passing Book~V legitimacy.  The two remaining
open obligations are O-LING.REC (recursion) and O-LING.SEM (semantic
grounding), neither of which is needed for the
contrast-compositionality endpoint.
\end{remark}

\subsection{O-LING.REC and O-LING.SEM: Remaining Frontiers}
\label{sec:ling-frontiers}

\begin{openobligation}[\status{O}]\label{ob:ling-REC-open}
\textup{(O-LING.REC: Recursion --- Remains Open.)}
The extension of the contrast-compositionality endpoint to infinite
generativity requires: (1) a certified iterated-assembly family
admitting self-embedding without bound; and (2) a Book~VI bridge
theorem licensing the infinite-generativity notion as the
faithful asymptotic limit of a certified sequence of finite assembly
operations.  No such bridge theorem is available in the current
canon.  O-LING.REC is not discharged here and is expected to remain
permanently open at the finite-observable level.
\end{openobligation}

\begin{openobligation}[\status{O}]\label{ob:ling-SEM-open}
\textup{(O-LING.SEM: Semantic Grounding --- Remains Open.)}
The extension of the contrast-compositionality endpoint to semantic
grounding requires: (1) a quotient-visible meaning assignment sourced
through $\mathbf{U}$; (2) no imported semantic domain, truth-value
space, or model-theoretic universe; and (3) a certified probe-response
structure on assembled objects supporting the meaning assignment.
No such structure is available in the current canon.  O-LING.SEM is
not discharged here.  It is the deepest downstream obligation in the
linguistics branch.
\end{openobligation}

%% ---------------------------------------------------------------
\section{SCC-Mediated Recursion}\label{sec:ling-REC-SCC}
%% ---------------------------------------------------------------

This section splits the monolithic recursion obligation O-LING.REC
into two parts and conditionally discharges the finite-stage SCC-carrier
part.  Strong infinite generativity retains its own named frontier.

\[
  \mathrm{O\text{-}LING.REC}
  \;=\;
  \mathrm{O\text{-}LING.REC}_{\mathrm{SCC}}
  \;+\;
  \mathrm{O\text{-}LING.REC}_{\infty}.
\]

\subsection{Iterated Assembly Chains}\label{sec:ling-REC-chains}

\begin{definition}[\status{D}]\label{def:ling-rec-chain}
\textup{(LING Finite Recursion Chain.)}
A \emph{LING finite recursion chain} is a sequence
$(X_n)_{n \ge 0}$ of certified LING objects in the
contrast-compositional regime $\mathcal{C}_{\mathrm{Ling}}^{cc}$
such that each transition $X_n \to X_{n+1}$ is generated by the
declared one-step compositional ledger $\mathcal{C}_{\mathrm{Ling}}$.
\end{definition}

\begin{definition}[\status{D}]\label{def:ling-rec-invariant}
\textup{(Transported Recursion Invariant.)}
The \emph{finite-stage recursion record} at depth $n$ is
\[
  T_n^{\mathrm{rec}} \;:=\;
  \bigl(
    [X_n]_K,\;
    \mathcal{C}_{\mathrm{Ling}}(X_n, X_{n+1}),\;
    \Delta_{\mathrm{Ling}}(X_n, X_{n+1}),\;
    \Gamma_{\mathrm{Ling}}|_{\le n}
  \bigr).
\]
This is not an infinite grammar.  It is a finite observable record at
depth $n$.  The \emph{transported recursion invariant} is the family
$T^{\mathrm{rec}} = (T_n^{\mathrm{rec}}, \rho_{n,n+1})_{n \ge 0}$,
where $\rho_{n,n+1} : T_n^{\mathrm{rec}} \to T_{n+1}^{\mathrm{rec}}$
is the declared transport map induced by one-step composition and
Book~III transport.
\end{definition}

\begin{theorem}[\status{C}]\label{thm:ling-rec-finite-stages-visible}
\textup{(Finite-Stage Visibility.)}
\textup{[Conditional on O-LING.CR through O-LING.END discharged.]}
Every $T_n^{\mathrm{rec}}$ is quotient-visible and source-descended.
No hidden alphabet, grammar, feature matrix, or semantic domain is
used.
\end{theorem}

\begin{proof}
Each $X_n$ is visible through the contrast kernel (O-LING.CR).  Each
step $X_n \to X_{n+1}$ is recorded by the compositional ledger
(O-LING.CL).  The contrast margin separates distinct certified classes
(O-LING.CM).  O-LING.NC blocks hidden parse rules.  O-LING.GD
identifies grammar descent from Book~III transport.  The LING endpoint
theorem already records that these ingredients are sufficient for the
contrast-compositional endpoint, which itself descends through Book~V
branch legitimacy.  Hence each finite recursion record is
source-descended.
\qed
\end{proof}

\begin{definition}[\status{D}]\label{def:ling-rec-support}
\textup{(SCC-Sufficient LING-Recursive Support.)}
A \emph{LING-recursive support family} $W^{\mathrm{rec}}$ consists
of the finite realized support components needed to recover $T_n^{\mathrm{rec}}$
for every finite stage $n$.  It is \emph{SCC-sufficient} when all
load-bearing support is exhausted by:
contrast outcome support,
composition-type support,
transport and defect-history support, and
ORA/CTM/TIN/ledger compatibility support.
This is the SCC support-sufficiency architecture specialized to LING
recursion; no semantic domain or hidden grammar is permitted.
\end{definition}

\begin{definition}[\status{D}]\label{def:ling-rec-carrier}
\textup{(SCC-Recursive LING Carrier.)}
A \emph{SCC-recursive LING carrier} $K^{\mathrm{rec}}$ is a lawful
faithful SCC carrier supporting $T^{\mathrm{rec}}$, satisfying:
finite-stage faithfulness, extension stability, no-hidden-grammar
discipline, and SCC support sufficiency.  It is \emph{minimal} when
no proper lawful subcarrier remains faithful.
\end{definition}

\begin{theorem}[\status{C}]\label{thm:ling-rec-support-sufficiency}
\textup{(LING-Recursive Support Sufficiency.)}
\textup{[Conditional on O-LING.END discharged; no hidden grammar,
hidden semantic domain, or undeclared compositional pathway.]}
The LING-recursive support family $W^{\mathrm{rec}}$ is SCC-sufficient.
\end{theorem}

\begin{proof}
Every load-bearing part of $T^{\mathrm{rec}}$ is one of four things:
a certified contrast outcome,
a certified one-step assembly type,
a transported defect or history record, or
a compatibility condition ensuring the update path is lawful.
Any additional support would be invisible in the contrast kernel,
compositional ledger, transport history, or compatibility ledger; it
would therefore be a hidden grammar, parse rule, feature system, or
semantic supplement.  That violates LING no-hidden-channel discipline
(O-LING.NC) and the SCC anti-smuggling rule.  Therefore the four
support types exhaust all load-bearing recursion support.
\qed
\end{proof}

\begin{theorem}[\status{C}]\label{thm:ling-rec-finite-localizability}
\textup{(LING-Recursive Finite Localizability.)}
\textup{[Conditional on Thm.~\ref{thm:ling-rec-support-sufficiency}
and the SCC finite-localizability chain (SCC Branch,
\texttt{thm:scc-sv}, \texttt{thm:scc-tl}, \texttt{thm:scc-fl}).]}
For every finite depth $n$, there exists a finite realized-support
subfamily $S_n \subseteq W^{\mathrm{rec}}$ such that $T_n^{\mathrm{rec}}$
is recoverable from $S_n$-level data alone.
\end{theorem}

\begin{proof}
At each depth $n$, the record $T_n^{\mathrm{rec}}$ contains finitely
many contrast classes, composition events, defect entries, and
compatibility checks.  The SCC branch packages the finite-localizability
chain as SCC-SV $\to$ SCC-TL $\to$ SCC-FL, using witness discreteness,
type compression, transport localization, and ledger discreteness.
Applying that chain to the LING-recursive invariant yields the finite
$S_n$ for each stage.
\qed
\end{proof}

\begin{theorem}[\status{C}]\label{thm:ling-rec-scc-carrier}
\textup{(SCC-Mediated LING Recursion.)}
\textup{[Conditional on: O-LING.CR through O-LING.END discharged;
Thms.~\ref{thm:ling-rec-support-sufficiency} and
\ref{thm:ling-rec-finite-localizability}; SCC-MCS available at
conditional force (\texttt{thm:scc-mcs}, SCC Branch); and
SCC-UC/source-descent verification in the applied SCC arena.]}
There exists a unique minimal SCC-recursive LING carrier
$M_{\mathrm{Ling-rec}}$, and every finite self-embedding stage
factors through it.  No hidden grammar or semantic domain is imported.
Therefore
\[
  \mathrm{O\text{-}LING.REC}_{\mathrm{SCC}} : [O] \to [C].
\]
\end{theorem}

\begin{proof}
Finite-stage visibility (Thm.~\ref{thm:ling-rec-finite-stages-visible})
gives lawful finite records.  The transported recursion invariant
(Def.~\ref{def:ling-rec-invariant}) gives a lawful update structure.
Support sufficiency (Thm.~\ref{thm:ling-rec-support-sufficiency}) blocks
hidden grammar.  Finite localizability
(Thm.~\ref{thm:ling-rec-finite-localizability}) supplies finite support
at every stage, enabling SCC chain-lower-bound arguments.  SCC-MCS
gives a unique minimal faithful carrier; SCC uniqueness blocks carrier
inflation.

The no-semantic-smuggling check: the data carried by
$M_{\mathrm{Ling-rec}}$ are contrast classes, composition events,
transport histories, and compatibility conditions.  None assigns
meanings, truth conditions, or reference.  If a semantic object were
required, it would be extra support not in the SCC-sufficient list, i.e.,
hidden supplementation, excluded by Thm.~\ref{thm:ling-rec-support-sufficiency}.

Therefore SCC-recursive LING certifies iterated contrast-visible
composition as stable finite-stage extensibility through a unique
minimal SCC carrier, not as an imported actually-infinite grammar.
\qed
\end{proof}

\begin{openobligation}[\status{O}]\label{ob:ling-rec-infty}
\textup{(O-LING.REC$_\infty$: Strong Infinite Generativity --- Split and Partially Discharged.)}
Strong infinite generativity is split into three sub-obligations:
\begin{enumerate}[label=\textup{(\arabic*)}]
  \item \emph{O-LING.REC$_{\mathrm{fam}}$}: construct a certified
    iterated assembly family --- conditionally discharged below;
  \item \emph{O-LING.REC$_{\mathrm{noncollapse}}$}: prove unbounded
    finite-stage distinctness does not collapse under the observational
    quotient --- conditionally discharged for declared noncollapsing families;
  \item \emph{O-LING.REC$_{\mathrm{bridge}}$}: prove the Book~VI
    asymptotic coherence bridge --- \emph{remains open}.
\end{enumerate}
\end{openobligation}

%% ---------------------------------------------------------------
\section{Strong Recursion Sub-Obligations}\label{sec:ling-REC-infty}
%% ---------------------------------------------------------------

\subsection{Finite Iterated Assembly Family}

\begin{definition}[\status{D}]\label{def:ling-certified-iterated-assembly-family}
\textup{(Certified Iterated Assembly Family.)}
A \emph{certified iterated assembly family} is a sequence
$\mathcal{A} = (X_0 \to X_1 \to X_2 \to \cdots)$
such that each transition $X_n \to X_{n+1}$ is produced by the
declared one-step compositional ledger $\mathcal{C}_{\mathrm{Ling}}$.
Each finite stage has a quotient-visible record $T_n^{\mathrm{rec}}$
(Def.~\ref{def:ling-rec-invariant}).  This imports no infinite grammar;
it is only a certified sequence of finite assembly records.
\end{definition}

\begin{theorem}[\status{C}]\label{thm:ling-rec-finite-iterated-assembly}
\textup{(Finite Iterated Assembly Legitimacy.)}
\textup{[Conditional on O-LING.CR, O-LING.CL, O-LING.CM, O-LING.NC,
and O-LING.GD discharged.]}
Every finite prefix $X_0 \to \cdots \to X_n$ of a certified iterated
assembly family is lawful LING branch data.
Therefore $\mathrm{O\text{-}LING.REC}_{\mathrm{fam}} : [O] \to [C]$.
\end{theorem}

\begin{proof}
Each step is a declared one-step assembly, so O-LING.CL supplies
additive ledger accounting.  Each object is visible through the
contrast kernel, so O-LING.CR and O-LING.CM supply quotient-visible
contrast separation.  O-LING.NC blocks hidden parse rules, alphabets,
or feature matrices.  O-LING.GD identifies grammar descent as the
restriction of Book~III transport to the contrast-visible sector.
The LING endpoint theorem already records that these ingredients
are sufficient for the contrast-compositional endpoint, which descends
through Book~V branch legitimacy.  Hence every finite assembly prefix
is source-descended.
\qed
\end{proof}

\subsection{Noncollapse and Unbounded Finite Generativity}

\begin{definition}[\status{D}]\label{def:ling-recursion-noncollapse}
\textup{(Recursion Noncollapse.)}
An iterated assembly family $\mathcal{A}$ is \emph{noncollapsing}
if for every depth $N$ there exist $m, n \ge N$ such that
$[X_m]_K \neq [X_n]_K$, or the finite-stage assembly histories
remain distinguishable by declared contrast/compositional probes.
In plain terms: the sequence does not merely cycle through finitely
many observationally identical states.
\end{definition}

\begin{theorem}[\status{C}]\label{thm:ling-rec-noncollapse-criterion}
\textup{(Noncollapse Yields Unbounded Finite Generativity.)}
\textup{[Conditional on Thm.~\ref{thm:ling-rec-finite-iterated-assembly}
and the assembly family being declared noncollapsing.]}
A noncollapsing certified iterated assembly family realizes unbounded
finite generativity at the contrast-compositional level.
Therefore $\mathrm{O\text{-}LING.REC}_{\mathrm{noncollapse}}$ is
conditionally discharged for declared noncollapsing families.
\end{theorem}

\begin{proof}
At every finite depth $N$, noncollapse supplies a later finite-stage
distinction visible to the declared probes.  Because contrast classes
are quotient-visible and distinct classes have positive contrast margin
(O-LING.CM), the family exhibits arbitrarily deep finite-stage novelty
without hidden grammar.  This proves \emph{unbounded finite generativity},
not actual infinite generativity.
\qed
\end{proof}

\subsection{Book~VI Asymptotic Bridge}

\begin{definition}[\status{D}]\label{def:ling-asymptotic-recursion-bridge}
\textup{(Asymptotic Recursion Bridge.)}
A \emph{Book~VI recursion bridge} for a noncollapsing certified
assembly family is a theorem showing that the sequence
$(T_n^{\mathrm{rec}})_{n \ge 0}$ is asymptotically coherent in
a declared normed recursion interface regime
$\mathcal{R}_{\mathrm{LingRec}}$, and admits an effective limit
object $T_\infty^{\mathrm{rec}}$ such that:
\begin{enumerate}[label=\textup{(\arabic*)}]
  \item every finite stage is recovered as a finite approximation;
  \item unbounded finite generativity is preserved in the limit;
  \item no infinite grammar is imported as primitive;
  \item the effective limit is shorthand for the asymptotic behavior
    of certified finite assembly.
\end{enumerate}
\end{definition}

\begin{theorem}[\status{B}]\label{thm:ling-rec-infty-bridge}
\textup{(Infinite Generativity Bridge.)}
\textup{[Bridge: conditional on a certified iterated assembly family
being noncollapsing and its finite-stage recursion records being
asymptotically coherent in the Book~VI sense in the declared regime
$\mathcal{R}_{\mathrm{LingRec}}$.]}
Strong infinite generativity is licensed as an effective asymptotic
object, not as a primitive grammar.
\end{theorem}

\begin{proof}
Book~VI licenses continuum or limit language only when it faithfully
encodes the asymptotic behavior of certified discrete structure in a
declared interface regime.  It explicitly says branch books must
verify the relevant asymptotic coherence and stabilization hypotheses
themselves.  Given those hypotheses, the infinite-generativity object
is not imported; it is the effective encoding of the certified
noncollapsing finite-stage family.
\qed
\end{proof}

\begin{openobligation}[\status{O}]\label{ob:ling-rec-infty-bridge}
\textup{(O-LING.REC$_{\mathrm{bridge}}$: Asymptotic Coherence --- Remains Open.)}
The minimal remaining recursion frontier is: prove that a noncollapsing
certified iterated LING assembly family satisfies Book~VI asymptotic
coherence in the declared recursion interface regime
$\mathcal{R}_{\mathrm{LingRec}}$.  This is the exact missing ingredient
for Thm.~\ref{thm:ling-rec-infty-bridge} to produce a fully licensed
infinite-generativity claim.
\end{openobligation}

\begin{remark}[\status{R}]\label{rem:ling-REC-full-status}
\textbf{Full recursion status after sub-split.}
The LING recursion picture is now three-tiered:
\begin{enumerate}[label=\textup{(\alph*)}]
  \item \emph{SCC-carrier recursion} (O-LING.REC$_{\mathrm{SCC}}$):
    conditionally discharged --- iterated composition admits a unique
    minimal SCC carrier; every finite stage factors through it;
    conditional on SCC-MCS.
  \item \emph{Noncollapsing finite generativity}
    (O-LING.REC$_{\mathrm{fam}}$ + O-LING.REC$_{\mathrm{noncollapse}}$):
    conditionally discharged for declared noncollapsing families ---
    certified finite iterated assembly is lawful; noncollapse implies
    unbounded finite novelty.
  \item \emph{Strong infinite generativity}
    (O-LING.REC$_{\mathrm{bridge}}$): the Book~VI asymptotic coherence
    bridge remains open.  This is the minimal remaining frontier for
    the full infinite-generativity claim.
\end{enumerate}
\end{remark}



%% ---------------------------------------------------------------
\section{Context-Response Semantic Grounding}\label{sec:ling-SEM-discharge}
%% ---------------------------------------------------------------

This section reduces O-LING.SEM from a vague ``find a meaning
assignment'' obligation to a precise conditional theorem with a named
remaining construction problem.  The key move: treat semantics as
a probe-response object, not as imported denotation.

\subsection{Semantic Probe Family}\label{sec:ling-SEM-probes}

\begin{definition}[\status{D}]\label{def:ling-context-frame}
\textup{(LING Context Frame.)}
A \emph{LING context frame} is a finite triple
$\mathcal{C} = (L, \Box, R)$
where $L$ and $R$ are finite assembled LING objects certified by the
compositional ledger, and $\Box$ is one declared insertion site.  For
an assembled object $X$, write $\mathcal{C}[X] = L \cdot X \cdot R$.
This uses only the certified compositional ledger (O-LING.CL).
\end{definition}

\begin{definition}[\status{D}]\label{def:ling-context-response-probe}
\textup{(Context-Response Probe.)}
A \emph{semantic context-response probe}
$P_{\mathcal{C},A}(X) := A\bigl([\mathcal{C}[X]]_K\bigr)$
is indexed by a context frame $\mathcal{C} = (L,\Box,R)$ and a
declared admissible contrast question $A$ on the assembled output.
The probe asks: which contrast class does $X$ produce when inserted
into a declared finite contrast-compositional context?  No truth
value, possible world, reference object, or model-theoretic domain
is used.
\end{definition}

\begin{definition}[\status{D}]\label{def:ling-sem-response-equivalence}
\textup{(Semantic Response Equivalence.)}
For assembled LING objects $X, Y$, define
\[
  X \sim_{\mathrm{sem}} Y
\]
iff $P_{\mathcal{C},A}(X) = P_{\mathcal{C},A}(Y)$ for every declared
context-response probe.  The \emph{semantic value} is
$\mu_{\mathrm{Ling}}(X) := [X]_{\sim_{\mathrm{sem}}}$.  This is
meaning as response-equivalence, not meaning as external reference.
\end{definition}

\subsection{Discharge on Declared Context-Response Regime}
\label{sec:ling-SEM-theorems}

\begin{theorem}[\status{C}]\label{thm:ling-sem-ctx-existence}
\textup{(Nontrivial Context-Probe Existence.)}
\textup{[Conditional on O-LING.CR and O-LING.CL discharged; the
admissible contrast probe family being nontrivially separating; and
the compositional ledger containing at least one nondegenerate context
$\mathcal{C}$ such that insertion can change the contrast class of the
assembled output.]}
The family $\mathcal{P}_{\mathrm{sem}}^{\mathrm{ctx}} :=
\{P_{\mathcal{C},A}\}$ is a nontrivial admissible semantic probe family.
\end{theorem}

\begin{proof}
The probe $P_{\mathcal{C},A}$ is a composition of two licensed
operations: (i) $X \mapsto L \cdot X \cdot R$ by the compositional
ledger theorem; (ii) $Y \mapsto A([Y]_K)$ by the contrast realization
theorem.  Nontriviality follows from the assumed nondegenerate context:
there exist $X, Y$ with $[\mathcal{C}[X]]_K \neq [\mathcal{C}[Y]]_K$.
By the contrast margin theorem, distinct contrast classes are positively
separated; hence at least one context-response probe separates $X$ and $Y$.
\qed
\end{proof}

\begin{theorem}[\status{C}]\label{thm:ling-sem-ctx-extension}
\textup{(Context-Extension Closure.)}
\textup{[Conditional on Thm.~\ref{thm:ling-sem-ctx-existence} and
the LING compositional ledger being closed under finite declared assembly.]}
The context-response probe family
$\mathcal{P}_{\mathrm{sem}}^{\mathrm{ctx}}$ is closed under
context extension: placing a context inside a larger certified context
yields another admissible context-response probe.
\end{theorem}

\begin{proof}
Let $\mathcal{C} = (L,\Box,R)$ be a certified context and $L', R'$
certified assembled objects.  By compositional ledger additivity,
$L' \cdot L$ and $R \cdot R'$ are certified finite assemblies; hence
$\mathcal{C}' = (L' \cdot L, \Box, R \cdot R')$ is a certified finite
context.  Since the probe family includes all finite certified contexts,
$P_{\mathcal{C}',A}$ belongs to $\mathcal{P}_{\mathrm{sem}}^{\mathrm{ctx}}$.
\qed
\end{proof}

\begin{theorem}[\status{C}]\label{thm:ling-sem-ctx-congruence}
\textup{(Semantic Congruence.)}
\textup{[Conditional on Thm.~\ref{thm:ling-sem-ctx-extension}.]}
Semantic response equivalence is stable under declared composition:
\[
  X \sim_{\mathrm{sem}} X',\quad Y \sim_{\mathrm{sem}} Y'
  \quad\Rightarrow\quad X \cdot Y \sim_{\mathrm{sem}} X' \cdot Y'.
\]
\end{theorem}

\begin{proof}
Apply any probe $P_{\mathcal{C},A}$ to $X \cdot Y$.  Then
$\mathcal{C}[X \cdot Y]$ is a context-extension probe on the pair
$(X,Y)$.  Since the family is closed under context extension
(Thm.~\ref{thm:ling-sem-ctx-extension}), the derived tests are
declared context-response probes.  Since $X \sim_{\mathrm{sem}} X'$
and $Y \sim_{\mathrm{sem}} Y'$, substitution preserves all declared
context responses.  Hence $P_{\mathcal{C},A}(X \cdot Y) =
P_{\mathcal{C},A}(X' \cdot Y')$ for every declared probe, giving
$X \cdot Y \sim_{\mathrm{sem}} X' \cdot Y'$.
\qed
\end{proof}

\begin{theorem}[\status{C}]\label{thm:ling-sem-source-descended}
\textup{(Source-Descended Semantics.)}
\textup{[Conditional on O-LING.END discharged and every semantic
probe acting only on certified contrast-compositional endpoint data.]}
$\mu_{\mathrm{Ling}}$ is source-descended through $\mathbf{U}$.
\end{theorem}

\begin{proof}
The contrast-compositional endpoint is already source-descended via
the LING branch-legitimacy chain.  A probe acting only on that
endpoint is downstream of the same source image.  Any semantic
distinction not detectable by declared probes would require an
undeclared semantic channel, excluded by the no-semantic-primitive
standing rule (Sr.~\ref{sr:ling-no-semantic-primitive}) and Book~V
hidden-supplement prohibition.
\qed
\end{proof}

\begin{theorem}[\status{C}]\label{thm:ling-sem-context-grounding}
\textup{(Context-Response Semantic Grounding.)}
\textup{[Conditional on O-LING.CR through O-LING.END discharged;
Thms.~\ref{thm:ling-sem-ctx-existence}, \ref{thm:ling-sem-ctx-extension},
\ref{thm:ling-sem-ctx-congruence}, \ref{thm:ling-sem-source-descended};
and no external semantic domain imported.]}
Define $\mu_{\mathrm{Ling}}(X) := [X]_{\sim_{\mathrm{sem}}}$.
Then $\mu_{\mathrm{Ling}}$ is a lawful semantic grounding object:
it is quotient-visible, source-descended through $\mathbf{U}$,
stable under declared composition, and introduces no external semantic
domain.  Therefore, relative to the declared nondegenerate
context-response regime:
\[
  \mathrm{O\text{-}LING.SEM}_{\mathrm{ctx}} : [O] \to [C].
\]
\end{theorem}

\begin{proof}
Thm.~\ref{thm:ling-sem-ctx-existence}: probes are admissible and
nontrivial.  Thm.~\ref{thm:ling-sem-ctx-extension}: family is
stable under context extension.  Thm.~\ref{thm:ling-sem-ctx-congruence}:
semantic equivalence is a congruence for declared composition.
Thm.~\ref{thm:ling-sem-source-descended}: $\mu_{\mathrm{Ling}}$
descends from $\mathbf{U}$.  No background semantic domain is
used: $\mu_{\mathrm{Ling}}(X)$ is the equivalence class of $X$
under declared context-response behavior.  Any semantic distinction
not detected by declared probes would require a hidden semantic
channel; that is excluded by Sr.~\ref{sr:ling-no-semantic-primitive}.
\qed
\end{proof}

\begin{remark}[\status{R}]\label{rem:ling-SEM-scope}
\textbf{Semantic scope note.}
Thm.~\ref{thm:ling-sem-context-grounding} gives the first lawful
semantic grounding layer: meaning as context-response equivalence.
It does not give reference, truth conditions, intentionality, or
model-theoretic semantics.  The semantic ladder is:
\[
  \text{contrast} \;\Rightarrow\;
  \text{composition} \;\Rightarrow\;
  \text{context-response meaning} \;\Rightarrow\;
  \underbrace{\text{reference-bearing probes}}_{\mathrm{O\text{-}LING.SEM}_{\mathrm{ref}}}
  \;\Rightarrow\;
  \underbrace{\text{truth-condition structure}}_{\mathrm{O\text{-}LING.SEM}_{\mathrm{truth}}}.
\]
The two downstream frontiers are open obligations.
\end{remark}

\begin{openobligation}[\status{O}]\label{ob:ling-SEM-ref-open}
\textup{(O-LING.SEM$_{\mathrm{ref}}$: Reference-Bearing Probes --- Remains Open.)}
Constructing a nontrivial reference-probe family that maps assembled
LING objects to stable selections of certified observable objects
within an already-certified arena, without importing an external
semantic domain, is a named open obligation.  This would upgrade
context-response semantics to reference-bearing semantics.
\end{openobligation}

%% ---------------------------------------------------------------
\section{Updated LING Closure Ledger}\label{sec:ling-closure-updated}
%% ---------------------------------------------------------------

\begin{remark}[\status{R}]\label{rem:ling-closure-updated}
\textbf{Updated discharge status after synthesis pass.}

\medskip
\noindent
\begin{tabular}{|l|l|p{7.5cm}|}
\hline
\textbf{Label} & \textbf{Status} & \textbf{Description} \\
\hline
O-LING.CR   & [C] & Contrast Realization \\
O-LING.CL   & [C] & Compositional Ledger \\
O-LING.CM   & [C] & Contrast Margin \\
O-LING.NC   & [C] & No Hidden Channel \\
O-LING.GD   & [C] & Grammar Descent \\
O-LING.END  & [C] & Contrast-compositionality endpoint \\
O-LING.REC$_{\mathrm{SCC}}$ & [C] & SCC-carrier recursion (conditional on SCC-MCS) \\
O-LING.REC$_{\mathrm{fam}}$ & [C] & Finite iterated assembly is lawful LING data \\
O-LING.REC$_{\mathrm{noncollapse}}$ & [C] & Noncollapsing families yield unbounded finite generativity \\
O-LING.REC$_{\mathrm{bridge}}$     & [O] & Book~VI asymptotic coherence bridge (remains open) \\
O-LING.SEM$_{\mathrm{ctx}}$ & [C] & Context-response semantic grounding \\
O-LING.REC$_\infty$       & [O] & Strong infinite generativity (Book~VI bridge needed) \\
O-LING.SEM$_{\mathrm{ref}}$ & [O] & Reference-bearing probes \\
O-LING.SEM$_{\mathrm{truth}}$ & [O] & Truth-condition structure \\
\hline
\end{tabular}

\medskip
The LING branch now has three tiers: (1) contrast-compositional
endpoint (conditionally closed); (2) SCC-carrier recursion
(conditionally discharged, gated on SCC-MCS) and context-response
semantics (conditionally discharged on the declared nondegenerate
context-response regime); (3) strong infinite generativity and
reference/truth semantics (genuinely open frontiers).
\end{remark}

\end{document}
